\section{Introduction}

\subsection{Overview}

VerletDrift is an interactive driving simulator built with vanilla HTML5 and JavaScript, implementing realistic vehicle dynamics through \textit{Verlet integration} and a sophisticated tire model based on the Pacejka magic formula. This document comprehensively documents the 43+ mathematical equations underlying the simulator, organized by functional category and cross-referenced with source code.

\subsection{Scope}

This reference covers:
\begin{itemize}
  \item \textbf{Core Physics}: Verlet particle integration, constraint solving, and rigid-body dynamics
  \item \textbf{Vehicle Systems}: Engine torque curves, transmission ratios, clutch engagement, and drivetrain
  \item \textbf{Tire Dynamics}: Pacejka magic formula, slip angles, grip saturation, and force relaxation
  \item \textbf{Steering}: Self-aligning torque (SAT), pneumatic trail, and steering column dynamics
  \item \textbf{Resistance Forces}: Rolling resistance, aerodynamic drag, and weight transfer
  \item \textbf{Camera \& Rendering}: Spring-damper systems, motion blur, and gauge needle animation
  \item \textbf{Scoring \& Effects}: Drift intensity metrics, particle systems, and game mechanics
\end{itemize}

\subsection{Unit System}

All equations use the \textbf{International System of Units (SI)}:
\begin{table}[h]
  \centering
  \begin{tabular}{ll}
    \toprule
    \textbf{Quantity} & \textbf{Unit} \\
    \midrule
    Distance & meter ($\unit{m}$) \\
    Time & second ($\unit{s}$) \\
    Mass & kilogram ($\unit{kg}$) \\
    Force & newton ($\unit{N} = \unit{kg \cdot m/s^2}$) \\
    Torque & newton-meter ($\unit{N \cdot m}$) \\
    Velocity & meter per second ($\unit{m/s}$) \\
    Angular velocity & radian per second ($\unit{rad/s}$) \\
    Acceleration & meter per second squared ($\unit{m/s^2}$) \\
    Moment of inertia & kilogram meter squared ($\unit{kg \cdot m^2}$) \\
    \bottomrule
  \end{tabular}
\end{table}

Speed in kilometers per hour (km/h) is converted internally: $v_{\unit{m/s}} = v_{\unit{km/h}} / 3.6$.

\subsection{Notation Conventions}

\subsubsection{Vectors and Scalars}
\begin{itemize}
  \item \textbf{Vectors}: Denoted as boldface, e.g., $\vv{v}$ for velocity vector
  \item \textbf{Scalars}: Italic for variables, e.g., $v$ for speed magnitude
  \item \textbf{Components}: Subscript notation, e.g., $v_x$, $v_y$, $v_{\text{long}}$, $v_{\text{lat}}$
\end{itemize}

\subsubsection{Time Derivatives}
\begin{itemize}
  \item $\dot{x} = \frac{\dd x}{\dd t}$ — first derivative (velocity)
  \item $\ddot{x}{t} = \frac{\dd^2 x}{\dd t^2}$ — second derivative (acceleration)
  \item $\dddot{x} = \frac{\dd^3 x}{\dd t^3}$ — third derivative (jerk)
\end{itemize}

\subsubsection{Common Abbreviations}
\begin{itemize}
  \item $\Delta$: Change in quantity (e.g., $\Delta v = v_{\text{new}} - v_{\text{old}}$)
  \item $N$: Normal load (force perpendicular to surface)
  \item $\alpha$: Slip angle (angle of tire relative to velocity direction)
  \item $\omega$: Angular velocity
  \item $I$: Moment of inertia
  \item $\tau$: Torque or time constant
  \item $\zeta$: Damping ratio (dimensionless)
  \item $\omega_0$: Natural frequency (rad/s)
\end{itemize}

\subsection{Code References}

Throughout this document, equations are cross-referenced to source code using the notation \coderef{file.js}{N}, indicating a specific line in the JavaScript source file. The complete codebase is available in the VerletDrift repository at \url{https://github.com/Quantum-Innovations/VerletDrift}.

\subsection{Structure of This Document}

Each equation section follows this structure:

\begin{enumerate}
  \item \textbf{Equation}: Mathematical formulation with proper notation
  \item \textbf{Physical Meaning}: Interpretation of what the equation represents
  \item \textbf{Source Code}: Code snippet showing the implementation
  \item \textbf{Parameters}: Definitions of all variables and constants
  \item \textbf{Notes}: Implementation details, approximations, or special cases
\end{enumerate}

\subsection{Verification}

All equations have been extracted and verified against:
\begin{itemize}
  \item Source code implementations in \texttt{physics.js}, \texttt{renderer.js}, \texttt{ui.js}, and supporting modules
  \item Mathematical consistency (dimensional analysis, limiting cases)
  \item Academic literature (where applicable) for standard models
\end{itemize}

\newpage
